\chapter{Як лагодити машину без схеми}
% full version: chapters/17_debugging.tex

\pencilfig{17}{\pencilart{17}}{Плата без схеми: щуп чує лише кілька проводів із мільярдів.}

Уявіть, що вам дали робочу плату без схеми й попросили її полагодити. Інженер має чотири інструменти: щуп, щоб послухати сигнал; генератор, щоб подати свій сигнал; можливість вимкнути деталь і подивитися, що зламається; і доступ до регістрів керування. Саме цими інструментами вивчають і мозок.

\section*{Щуп на тисячу нейронів}

Електрод, устромлений у мозок, чує клацання кількох нейронів поруч. Здавалося б, на тлі вісімдесяти шести мільярдів тисяча електродів --- ніщо. Але їх вистачає, щоб керувати курсором на екрані. Чому? Сусідні нейрони шумлять схоже, і активність, що керує рукою, насправді описується лише одним-двома десятками спільних змінних. Не треба слухати всіх, досить почути ці кілька «голосів».

Сирий сигнал із тисячі електродів --- це сотні мільйонів бітів на секунду: радіоканалові з-під черепа такого не пропустити. А самі клацання несуть приблизно в тисячу разів менше. Тому клацання вигідно виокремлювати просто на імплантаті й відсилати назовні лише їх --- той самий прийом стиснення, що й в оці.

\section*{Що робить Neuralink}

Компанія Neuralink узялася за інженерний бік задачі: тонкі гнучкі нитки замість жорстких голок, робот, який їх вживлює, і герметичний бездротовий імплантат. У її опублікованій праці описано систему, з'єднану з комп'ютером кабелем, а стиснення й бездротовий зв'язок названо планами. Дані про людину з паралізованими руками, яка керує курсором, поки відомі здебільшого зі звітів самої компанії; за ними в імплантаті близько тисячі каналів на кількох десятках ниток, і частина ниток після вживлення змістилася. Сама ідея не нова: академічні групи з жорсткими ґратками приблизно із сотні електродів уже давали людям керувати курсором, «писати» літери силою думки й перетворювати спроби говорити на текст на екрані зі швидкістю близько шістдесяти слів на хвилину.

\section*{Читання і запис}

Декодер читає частоти клацань багатьох нейронів і отримує менше ніж десять бітів на секунду --- мізерну частку того, що несуть самі нейрони. Зате і мозок, і декодер учаться одне в одного: людина пристосовується до декодера, декодер --- до людини. Два учні, що підлаштовуються один під одного, --- тонка задача: навіть коли кожен окремо стійкий, разом вони можуть розгойдатися.

Записувати в мозок важче, ніж читати. Струм через електрод збуджує все навколо, а не один потрібний нейрон. Можна «намалювати» точки світла, як пікселі, що світяться, і людина розрізнить прості літери. Але записати стиснений код ока так, щоб мозок побачив картинку, поки не вміє ніхто.

\section*{Чого щуп не бачить}

Електрод чує телефонну мережу даних. Але радіостанції --- дофамін, серотонін, норадреналін, ацетилхолін --- це концентрації речовин, їх щуп не чує. Не бачить він і сили зв'язків, де зберігається вся пам'ять. І не бачить болю так, як його відчуває людина. Налагоджувач, який бачить дані, але не бачить регістрів, завжди дивуватиметься, чому той самий вхід дає різні відповіді.

\fullversion{17}
